% 原书第 161 页（PDF 第 174 页）
\BookPage{161}{174}
\BookHeading{清代籍没律例与抄家案件}{清代籍没律例与抄家案件}{柏桦*}
抄家是中国古代长期存在的刑罚。大体上是王朝在对一些刑事犯罪，除了判处笞、杖、徒、流、死等刑罚以外，往往附带处以一定的罚金，作为给受害人的物质赔偿（如常见的烧埋费、汤药费等），或直接没入官库。但除了对大逆、谋反等十恶不赦的人犯以外，是极少探用抄家办法的。对什么人和犯什么罪应予抄家惩处，《大清律例》并无明确的条款。从分散在中国第一历史档案馆、台湾故宫博物院图书文献处的《上谕档》、《内务府奏销档》、《军机处朱批奏折·内政职官类》、《法律类》、《六科史书》等档案中，可以看到皇帝宣布某些官僚贵族的罪行，决定治罪并抄家的谕旨；奉派前往执行查抄任务的大臣或地方官奏报查抄经过；还有查抄财产的清单和对这些财产的处置等。处理方式有些是依据《大清律例》，有些则出于「圣裁」，处置过程所体现的政治与经济因素，值得思考。

抄家也称为「籍没」、「抄没」，其特点是加大惩罚力度，不仅是以罪犯本人为对象，而是扩大到罪犯的家庭，乃至于家族。抄家也不仅仅以财产为对象，还包括犯罪者家庭或家族人口的命运。凡是被判决抄家者，其决定权并不在刑部、三法司、督抚等主管司法的部门和人员，要取自上裁，由皇帝掌握最终的裁决权。

清代对重大犯罪实施抄家处罚是比较普遍的，其中对官员的抄家更是超过以往的朝代。这是因为满族在国家形成的过程中，上层统治者有相互兼并财产部属的习惯，在王朝建立过程中又与中国历代王朝皇权无限的传统结合在一起。因此可以说对官僚贵族的抄家既是当时权力和财产再分配的形式之一，也是官僚贵族们盛衰沉

\vfill
\begin{BookNotes}
\BookNote{*}{作者任职于南开大学法学院。}
\end{BookNotes}

% 原书第 162 页（PDF 第 175 页）
\BookPage{162}{175}
\noindent 浮的标志之一。

早在入关前后，满洲皇族内部，即在努尔哈齐诸子兄弟叔侄褚英、代善、莽古尔泰、皇太极、阿敏之间；其后在多尔衮、阿济格、多铎、豪格、索尼、济尔哈朗、硕托、阿达理之间，一直就存在复杂尖锐的矛盾，互相倾轧几成习俗。褚英、阿敏、索尼等获罪后，原属他们统率的军队，所拥有的金银财宝、牲畜、奴隶，都被继起的掌权人所占有。如多尔衮最盛的时候，除了直接统率正白、镶白两旗外，还严密控制住正黄、镶黄两旗，当时称为「节制四旗」，即掌握着满洲八旗军力的一半，而且都是精锐。但到他死后被削爵追究罪责时，除了正黄、正白、镶黄三旗军队被收归皇帝直辖，成为「上三旗」外，多尔衮及其弟阿济格，还有何洛会、胡锡、刚林、谭泰等亲信都被抄家。顺治帝在查办多尔衮的上谕中讲：「其所用仪仗音乐及卫从之人，俱僭拟至尊。盖造府第亦与宫阙无异，府库之财任意縻费，擅用织造缎匹，库贮银两珍宝不可胜计。」不但指斥其生前的仪仗、住房、穿带、用具都超过规定，属于「僭用」，而且将多尔衮当政时，把皇上侍臣伊尔登、陈泰一族及所属牛录收归该旗下也当为罪状，于是要以其人之道还治其人之身，所以要抄家，而许多亲信也被籍没家产，有些还将他们的妻子发配为奴。当年多尔衮主持抄了索尼、阿敏等人的家，如今顺治帝又抄了多尔衮等人的家。抄人家又被人抄家，成为满族统治集团在政争中的胜败，以及权力交替的明显标志。

在历史档案中，康熙、雍正、乾隆、嘉庆、道光等时期的抄家档案存留较多，其间发生过许多重大的抄家案件，仅据笔者查阅及抄录的300余件抄家档案，再结合清代的律例，将这些抄家案件粗分为五类：

第一，皇帝对当时最有权势的大臣探取断然措施、逮捕处死并予以抄家的。这里当属康熙八年（1669）逮捕并永远囚禁原首席辅政大臣鳌拜；雍正三年（1725）赐大将军年羹尧自尽；嘉庆四年（1799）处死首席大学士和珅等最为著名。他们被处置和抄家，都发生在新皇帝上台不久，正在权力交替和巩固权位的过程中，显然是出于从政治上加强皇帝专制权力的考虑。值得注意的是，这三个人被处罚并抄家，并不是以政治上的罪名，而是按照经济案处置的。

\vfill
\begin{BookNotes}
\BookNote{1.}{《清世祖实录》卷53，顺治八年二月己亥条。}
\end{BookNotes}

% 原书第 163 页（PDF 第 176 页）
\BookPage{163}{176}
第二，由于违反法纪而又失去皇帝信任的官员。不论是中央的内阁、军机处，还是部、院、寺、监、府的大臣或宗亲贵族，抑或对地方上省、司、道、府、州、县官吏，一旦有违法乱纪行为，往往都籍没抄家，而督抚等大员的抄家单往往是几次上报。我们在收集整理清代嘉庆以前，被朝廷惩处的中央和地方一二品满汉军政大员132人共109件经济犯罪案件时，发现除年羹尧（92款重罪）、蔡珽（18款重罪）、隆科多（41款重罪）、延信（20款重罪）、和珅（20款重罪）等少数人，被认为是因为政治的原因而牵涉到经济犯罪外，其他均被认为是受贿、勒派、克扣、挪用、侵贪等经济案件。这些案件涉及司法审判、行政管理、教育科举、经济活动、军事用兵等领域。132名大员中，地方官员有109人，占82.6\%，督抚占大多数。在抄家单上查到的督抚有：康熙时云贵总督蔡毓荣；雍正时直隶总督李维钧、甘肃巡抚胡期恒、四川巡抚蔡珽；乾隆时浙江巡抚卢焯（后被赦免，再任陕西巡抚时被二次抄家）、江南河道总督周学健、署理湖广总督塞楞、陕甘总督永常、甘肃巡抚鄂昌、云贵总督张广泗、河东河道总督李亨特、直隶总督杨灏、湖南巡抚蒋炳、江苏巡抚庄有恭、广西巡抚钱度、云贵总督李侍尧（后被赦免，再任两广总督时被二次抄家）、直隶总督杨景素、浙江巡抚王亶望和陈辉祖、四川总督勒尔谨等；嘉庆时云贵总督的鄂辉及富纲、陕西巡抚秦承恩、河南巡抚毕沅、浙江巡抚琅玕、湖北巡抚常明等。高级官员犯罪的原因非常复杂，被惩处的理由也不尽相同。一方面是「贿赂公行，已成积习」2，有主观和偶然的因素；另一方面是专制君主独揽大权，但也不能不赋予官吏以一些权力。「一种权力不过是国家的一般意志，另一种权力不过是这种意志的执行而已。」君主授予官吏的权力是执行自己的意志，但官吏拥有权力，就可以凭借权力攫取自己的利益，于是便有「属官情愿送礼」，亏空库银或因公挪用，军需紧急「并非侵蚀」4等非正常情况发生。如果忤逆于君主和上司，除了依律定罪，还要「从重归结」5。在驾驭与反驾驭的过程中，实际上是各种政治势力的较量，而不是认真履行法律。

\vfill
\begin{BookNotes}
\BookNote{2.}{（清）饶玉成编，《皇朝经世文编续集》卷19〈吏政·考察〉，载骆秉章四川吏治营务废弛已极疏，清光绪八年（1882)江右饶氏双峰书屋刊本。}
\BookNote{3.}{（法）孟德斯鸠着，张雁深译，《论法的精神》（上册），商务印书馆，1961年，第157页。}
\BookNote{4.}{《清圣祖实录》卷198，康熙三十九年三月丙申条。}
\BookNote{5.}{《清圣祖实录》卷251，康熙五十一年九月庚子条。}
\end{BookNotes}

% 原书第 164 页（PDF 第 177 页）
\BookPage{164}{177}
第三，有关「文字狱」的抄家案件。「文字狱」不是清代特有的，但对「文字狱」涉案者探取广泛株连和籍没抄家，可以说是清代的特色。康熙时期的庄廷锈「明史案」，杀了70多人，连刻书、买书的人也不放过；戴名世《子遗录》的「大逆案」，逮捕200多人，连早已死去的方孝标也被剖棺戮尸。雍正时期的汪景祺文章讥讽案、钱名世名教罪人案、查嗣庭科场试题案、陆生楠《通鉴论》案、谢济世注《大学》案、曾静和张熙「谋反」案等，都是著名的「文字狱」大案。乾隆时期的《字贯》、著作、家谱诸案迭生，而「伪孙嘉淦奏稿案」、「徐述夔一柱楼诗案」、「胡中藻坚磨生诗钞」等案件更扩大打击面。这祖孙三代大搞文字狱，在社会上造成很大的恐怖。这类案件查抄重点是图书、著作、诗文稿、来往书信等，对其他财产是先封锢，如乾隆帝批复直隶总督袁守侗报告查抄尹嘉铨原籍住宅经过的奏折中，对其财产要求「此皆不必动」6。从原故宫博物院文献馆20世纪30年代编的《清代文字狱档》来看，几乎每案中都有「严查有无违碍应缴书籍」；「细查有无不法字迹，片纸只字不得遗留」；「即墙壁窟穴中亦必详检无遗，倘致透漏风声，伊家得以藏匿，惟尔等是问」等批示。「文字狱」抄家重点虽然不在财产，但涉案者的家族，乃至宗族人口被籍没，或流放，或发配为奴，比一般抄没财产更凶残。

第四，因为各级官府仓库亏空钱粮数额巨大，责任者不能依限还清亏欠钱粮而被抄家的。清代府库出现亏空，先由亏空官本身追取，因此要先将亏空官任所及原籍财产查封，勒令完补，如果亏空官家产尽绝，着落保题之上司官均行赔补，所以在亏空案中被抄家的很多。在这类抄家案中，比较引人注目的是内务府人员。内务府是为皇帝、皇族、宫廷的生活服务机构，直属于皇帝，其下辖有：广储司、都虞司、掌仪司、会计司、庆丰司、慎刑司、营造司等七司，上驷院、奉宸苑、武备院等三院，此外还有三织造处（江宁、苏州、杭州）、三旗参领处、掌关防处、三旗庄头处、养心殿造办处、武英殿修书处、御茶膳房、升平署、御药房、咸安宫官学等诸多附属机构，而分置在各地的盐政、织造、海关监督等，大部分是由内务府人员来担任的。盐政、织造、海关是第一等肥缺，担任这类职务的官员，不但要认真揣摩迎合皇帝的喜好和需要，随时贡献各种奇珍异物及巨额金银，还必须给予中央

\vfill
\begin{BookNotes}
\BookNote{6.}{《殊批奏折》内政类，〈职官〉，乾隆四十六年三月二十五日（1781年4月18日），袁守侗奏。}
\end{BookNotes}

% 原书第 165 页（PDF 第 178 页）
\BookPage{165}{178}
\noindent 和地方上有权势的部门和人物一定的分润，本身更是奢侈腐败。这些官员及其家族骄奢淫逸的生活方式，曾经受到时人的艳羡，但到头来亏空公款的额度愈来愈大。这类在被抄家的内务府差委人员，少者亏欠「内帑」数万两，多者高达数百万两，显然是不能够满足各方面的需索，对皇帝也贡献有缺，于是被抄了家，以家财赔补历年积欠，与《红楼梦》作者曹雪芹有关的曹寅和李煦就是这样的下场。当然，因为亏空而被抄家顶债，不仅仅是经济上的原因，其中还有明显的政治原因，但引起曹、李两家被抄家清产的直接原因还是存在着巨额的亏空。《内务府奏销档》所载被抄家的，如雍正时粤海关监督祖秉圭；乾隆时苏州织造安宁，两淮盐政伊龄阿，管理荆关税务丽柱，淮安关监督伊拉齐，杭州织造隆升，粤海关监督郑伍赛、李永标、德魁；嘉庆时淮关监督及长芦盐政李如枚等，都被撤职、判刑和抄家。他们之中固然不乏巨贪，但也有被誉为「操行洁白，吏役欺罔，时加戒饬。捐廉优恤，以示鼓舞」的廉吏如淮安关监督伊拉齐等。翻阅这些曾经豪奢一时者的抄家清单，再与史料对照，发现他们财产与开销根本就不成正比。如乾隆四十五年（1780），在南巡至扬州时，行宫设在天宁寺，仅二月十五日的一顿饭，两淮盐政伊龄阿就提供了余馍、安膳桌、银葵花盒小菜、银碟小菜、粳米干膳、额食、奶子、领炉食、盘肉、羊肉四方等15桌菜肴；此外前来觐见的督抚以下官员还有20桌。号称不扰民的乾隆帝，没有吃伊龄阿提供的饭菜，为左右人亲自烹饪，其菜品有燕窝肥鸡丝热锅、火熏东坡鸭子、鹿筋酒烧鸡冠肉、羊肉片、炒苏蛋、春笋炒肉、蒸烧肥鸡羊乌叉、象眼小馒首、白面丝糕糜子米面糕等，这些材料当然要两淮盐政伊龄阿提供了8。此一餐真不知要花费多少中人之产，伊龄阿又如何不亏空？最终结果，不但这些官员难免抄家，就是与他们打交道的「皇商」、「盐商」、「洋商」，也不能够幸免于抄家籍没。

第五，因为谋反、叛逆、暴乱等反抗朝廷的团伙，罪大恶极而手段残忍的杀人犯。业师韦庆远先生在20世纪70年代主编的《康雍干时期城乡人民反抗斗争资料》中，对这一时期的农民、奴婢、手工业者、城镇商民反抗地主、奴隶主、

\vfill
\begin{BookNotes}
\BookNote{7.}{冒广生，《淮关小志》，方志出版社，2006年，第556页。}
\BookNote{8.}{参见中国第一历史档案馆、扬州市档案局（馆)，《清宫扬州御档选编》所载〈乾隆帝在天甯寺高旻寺行宫用膳形事，广陵书社，2009年，第193-194页。}
\end{BookNotes}

% 原书第 166 页（PDF 第 179 页）
\BookPage{166}{179}
\noindent 雇主、奸商、贪官等资料进行分类汇编。清王朝在处置这些案件时，往往以「谋反」罪名实行严厉打击，抄没家产和人口也是常见的。第一历史档案馆藏《刑科题本》，也有许多「谋反」、「谋逆」的案件，大多数被籍没财产人口。这类涉案人员因为财产无多，所以在田产入官以后，人口多被发遣或发配为奴。此外，在平定三藩、准噶尔贵族叛乱、大小金川等战役过程中，对叛乱者的财产人口也实行抄没。

把抄家档案分为五大类，实际上也不能够反映出抄家案件的全部内容，更不能够用来分析其中复杂的案情，也很少能够看到依据什么法律裁断的。因此，必须要知道清代的法规体系，了解案件发生的政治与社会背景。

二

清代律例规定的抄没刑罚，主要针对威胁到王朝统治的罪犯，而对于犯有重大罪行的各级官吏和贵族，与官府有密切关系的皇商、官商人等，以及文字狱中的「思想犯」，反抗王朝统治和破坏王朝统治秩序的反叛者等，也有抄没的规定。这些罪犯在刑罚之外都有籍没抄家。当然，也有在定案判罪之前即谕令籍没抄家者，这乃是搜获罪证和扩大株连的必要手段。

清顺治三年（1646）奏准，四年（1647）颁行「顺治律」以后，经过七十余年，于雍正元年（1723）开馆修律，三年（1725）颁行《大清律集解附例》，勒定清律436条的规模。乾隆元年（1736)再度开馆修律，五年（1740）颁行《大清律例》，此后律文基本不改，一直延续到清末修律。

在律文中，有关抄家在「给没赃物」、「隐瞒入官财产」、「谋反大逆」等条有规定。如「给没赃物」条，对贪赃枉法、私藏违禁物品、恐吓和诈欺取财、科敛和求索等行为，要将赃物及违禁物品籍没入官；对谋反、叛逆者不仅仅籍没财产，还要家口缘坐；对一般犯罪给别人造成损失者，追罚钱财，但不籍没财产。凡是籍没的财产，官物还官，私物还主。为此，律文还规定估赃的标准，一般物品按中等物价估定，雇工按一人一日的工钱计算，牛、马、驼、骡、车、船、碾、磨、店舍等按雇赁价值计算。「隐瞒入官财产」条规定，谋反、谋叛、奸党等在「十恶」的


% 原书第 167 页（PDF 第 180 页）
\BookPage{167}{180}
\noindent 罪犯，要抄没人口、财产，如果在抄没时有隐瞒财产和人口的行为，要予以治罪；对同情隐瞒的里长，知情的官吏也要治罪；对举报的人，如果举报得实有赏，举报不实则有罚，不举报则处以笞五十以下的刑罚。「谋反大逆」条，除谋反正犯凌迟处死之外，其家族、宗族16岁以上的男子皆斩，15岁以下的男子及所有女子（嫁人除外）都给付功臣之家为奴，财产入官。「谋叛」条，谋叛者不分首从皆斩，本家人口给付功臣之家为奴，财产入官，家族则流放。

康熙九年（1670），另行修订「刑部条例」，别自成《现行则例》，收有例文290条，从《大清律朱注广汇全书》及《古今图书集成·祥刑典》来看，则是将当时458条律文分别加以纂注，并在其后增加了比附条例。康熙十八年（1679），颁布上谕指出:「于定律之外，附设条例。」要求在定律之外，所有条例应去应存，着九卿詹事科道会同详细加酌定确议。经遵旨议定，在《大清律集解附例》之外另定《刑部见行则例》，在康熙二十九年（1690）再次刊行时，则名为《刑部新定现行则例》。根据康熙的指示，「命图纳、张玉书等为总裁，将律文则例，删定改正，汇为一书，而《刑部见行则例》之撤销，即始于此，自后终有清一代，各部署皆有则例，刑部但有《通行则例》与《章程》。不复以则例见称矣」9。乾隆五年（1740）以后，形成五年小修，十年大修的修例制度。律与例并行的用意是：「律为一定不易之成法，例为因时制宜之良规。故凡律所不备，必藉有例，以权其大小轻重之衡，使之纤悉比附，归于至当。」10清代条例庞杂而繁琐，每次修例皆以前次为基础，进行修改、修并、移改、续纂、删除等技术手段处理后11，例的数量和内容随时都处在动态的变化之中。例在法律上处于优先的地位，这是清代法律的一个特点，此点非常重要，不可忽视。研究清代法律不能忽视条例，否则不但不了解法律的具体运用，还会发生错误2。

\vfill
\begin{BookNotes}
\BookNote{9.}{王钟翰，《清史论集·清代则例及其政法关系之研究》，中华书局，2004年，第3册第1748页。}
\BookNote{10.}{乾隆元年傅鼐《奏请修定律例疏》，载田涛、郑秦点校，《大清律例》，法律出版社，1999年，第13页。}
\BookNote{11.}{修改，亦即将原来的例加以审视，对其条文略加修正；修并，就是对原来的例加以整合，将其中二条以上的合为一条；移改，也是对原来的例加以整合，将其中的一些条文移易其类属位置；续纂，这是指在原来的例中加入所没有的新内容；删除，将原来的例中的某些内容删削除去。}
\BookNote{12.}{参见瞿同祖，《瞿同祖法学论着集》，中国政法大学出版社，1998年，第423页。}
\end{BookNotes}

% 原书第 168 页（PDF 第 181 页）
\BookPage{168}{181}
涉及籍没抄家的条例很多，仅据上列律文所附条例来看，「给没赃物」经过修改、修并、移改、续纂、删除者就有34条，「隐瞒入官财产」有9条，「谋反大逆」有10条，「谋叛」有20条。这些条例在很大程度上补充律文之不足。如「谋反大逆」的条例一，对反逆案内干连流犯的无子之妻免流；条例二，对反逆案缘坐男犯15岁以下者送交内务府阉割，10岁以下则监禁到11岁再交内务府阉割；条例三，明确了缘坐16岁以上男犯发配为奴，妇女及男10岁以下给官为奴，男11-15岁则监禁到16岁，再发配为奴，而不再将他们阉割；条例四，反逆案内凌迟之犯的子孙交内务府阉割后发配为奴，10岁以下者监禁到11岁再阉割发配为奴，缘坐妇女发各省驻防给官员兵丁为奴。条例五，是比照反逆案的亲属，不阉割为奴，均改为流放；条例六，比照大逆定罪的罪犯不知情之祖父、父母无罪释放；条例七，比照反逆定罪的罪犯，其不知情尊长免缘坐；条例八，被怀挟私嫌诬陷者，不能够以缘坐律问拟；条例九，凡是比照反逆及谋叛定罪的罪犯家属一概免其缘坐；条例十，已经给功臣为奴的人犯，主人不能养赡，则发配为奴13。从这10条来看，与律文相比，多数是属于从轻处置，也有从重处置，而从例不断修订的情况来看，「条例则多从宽典」14，在逐渐减少缘坐的情况下，缩小了打击面。

清代律例适于使用的同时，与律例相关的各部、院、寺、监、府的「则例」是当时的行政准则，在政务具体执行过程中起到重要的作用。其中与抄家籍没财产人口有关的，当属《吏部处分则例》、《督捕则例》、《理藩院则例》中的一些规定。则例侧重于行政处分，即以罚俸、降级、革职为处分，如果则例中规定「议罪」，这就是行政处罚之外还有刑事处罚。由于刑事处罚与行政处罚分属不同的部门，在具体实施过程中，还有各部门之间协调的问题。比如说《吏部处分例·亏空分赔》规定，一旦出现府库亏空，就要将本犯拟罪监追，勒限完补，在勒限期间，要先将其家财产查封，如期不能够赔补足额，再查封其家族、宗族财产，到期这些财产不能抵充亏空，则要各上司赔补。再如《督捕则例》是在众多则例中属于影响较大的，其与律例虽然侧重方面有所不同，但更重要的是彼此相互关联融通，例

\vfill
\begin{BookNotes}
\BookNote{13.}{参见（清）官修，光绪《清会典事例》卷779〈刑部·刑律贼盗·谋反大逆〉，台湾新文丰出版公司据清光绪二十五年（1899）原刻本影印本。}
\BookNote{14.}{（清）薛允升，《读例存疑》卷25〈贼盗·谋反大逆〉，光绪三十一年（1905）京师刊本。}
\end{BookNotes}

% 原书第 169 页（PDF 第 182 页）
\BookPage{169}{182}
\noindent 如「窝逃及邻佑人等分别治罪」条，自顺治五年（1648）至顺治十四年（1657）修订9次，是专门针对窝家而言，从最初「窝家正法，妻子家产籍没给主，仍给一分与出首之人，邻右十家长等，各责四十板，流徙边远」的严厉刑罚，逐渐调整为不株连分家之父子兄弟，窝家免死板责同妻子一并流徙，十家长减责为二十板，窝逃人的家产给予逃人之主，房地入官，两邻、十家长也要受到板责。顺治十一年（1654）又开始重惩窝家，不但要将窝家的家属人口充发盛京，还对两邻、十家长、该管官施行罚银给逃人之主及出首之人，窝隐逃人本犯正法，家产房地入官，两邻各责四十板流徙，十家长责四十板，所罚之银入官；旗下人窝隐逃人，不但要鞭一百，还要罚银。顺治十三年（1656），十家长也要责四十板流徙，而顺治十四年（1657），窝犯免死板责刺字，家产人口给予八旗穷兵，并强调地方官、地邻的责任15。9次修订既反映缉捕逃人宽严变化的过程，又反映统治者对待窝家的态度，不将窝家正法，而没收窝家的财产，窝家及相关责任人的流徙充发，使八旗奴仆的来源增加。另如《理藩院则例》中规定在蒙古地区犯偷窃、强盗、发冢、探参、捕猎、杀人、放火等罪时，妻子畜产都将被籍没，不是给予事主，便是给予台吉，人口都被罚为奴，而当清王朝对蒙古地区控制力加强的时候，籍没财产给事主或台吉，人口则发遣西南或内地当差。凡此，都可以看出清代法规的多民族色彩，以及政治与法律制度所表现出的张驰有度、刚柔相济的特点。

事例的形成源起于律例所涉及的事和案。日本织田万认为：「清国成文法乃既饶多，不文法之势力亦且极其大矣。」16通过颁行事例，以画归一，以儆官邪。事例的形成有一定的程式，最终由君主颁行谕旨或批准遵行，充分体现着皇权至高无上的权威，体现着君主专制的高度集权。关于抄家籍没财产人口的事例也很多。例如顺治十五年（1658）定:「镶黄等三旗有籍没家产者，交与内该管衙门。」顺治十八年（1661）覆准:「除五旗重犯籍没家产，俱交各该都统、副都统拨给本旗外，其三旗重犯籍没家产，免交内该管衙门，交与该都统、副都统等请旨拨给。康熙七年（1668）覆准:「凡犯军机籍没家产者，除婢妾外，仍给予人口三双，牛

\vfill
\begin{BookNotes}
\BookNote{15.}{（清）官修，光绪《大清会典事例》卷855〈刑部·督捕例〉，台湾新文丰出版公司据清光绪二十五年（1899）原刻本影印本。}
\BookNote{16.}{（日）织田万着，李秀清、王沛点校，《清国行政法》，中国政法大学出版社，2003年，第48页。}
\end{BookNotes}

% 原书第 170 页（PDF 第 183 页）
\BookPage{170}{183}
\noindent 马各三匹并器械等件。17这反映出时代的特点，早期全部籍没而不留任何财产，后来给予一定的生活物品。再如，顺治时整饬吏治，凡是隐匿欺蒙钱粮款项，全部杀无赦，还「行籍没人口入官，家产充告发之赏」18。贪官「受赃至十两以上者，不分枉法、不枉法，俱籍没家产入官」19。另如，康熙初年，对喇嘛教实行严格管理，凡是喇嘛不领度牒，私收人为班第（喇嘛最低级，相当于徒弟），便将「喇嘛处绞，家产籍没」20。凡此，都说明事例的时效性，在位的皇帝颁布事例以后，其效力高于所有的法律，但随着时间的推移，统治者虽然没有明令废除该事例，但已经很少或根本不再援引了。

清代的省例是以地方性事务为规范对象，以地方行政性法规为主体，兼含少量地区性特别法的一种法规汇编，在各地司法、行政过程中具有重要作用。「省例则因地制宜，故凡一州一县议详之事件，足为程式者，皆探列，颁之全省，通行各属」21。因此「省例的设置，其首要的目的是因地制宜，结合各地方的不同情况作出便于执行的规定」22。省例作为朝廷法律的实施细则和补充规范，总体上与朝廷法律保持一致。但各省例在不悖于朝廷基本精神的前提下，也存在不少差异。

章程是律例之外的「通行」，补律例之所未尽，其特点是详于案而略于例。仅以《刑部通行章程》而言，就载有1838-1892年通行的行催汇题、窃毁电线、严禁非刑、酌议会匪、变通叩暗、监禁、扒窃、阵亡、军流脱逃、禁止滥用非刑、禁止吕宋赌票、京控、蠹役正法、军流徒犯脱逃解役、禁止待质公所、停解军遣、咨题案件咨覆限期、私铸银钱、严定武官扣饷缺额、缉捕盗贼、严禁非刑、盗墓未得财、整顿刑名、捉人勒赎之犯拟入情实、就地正法等百余种章程。王钟翰先生认为章程编纂性质不外七种:「一曰䌷绎例意；二曰申明定章；三曰比例不能画一而折

\vfill
\begin{BookNotes}
\BookNote{17.}{（清）官修，光绪《清会典事例》卷727〈刑部·名例·犯罪免发遣〉，台湾新文丰出版公司据清光绪二十五年（1899）原刻本影印本。}
\BookNote{18.}{《清世祖实录》卷48，顺治七年四月壬子条。}
\BookNote{19.}{《清世祖实录》卷95，顺治十二年十一月丁亥条。}
\BookNote{20.}{（清）官修，光绪《清会典事例》卷752〈刑部·户役·私创庵院及私度僧道〉，台湾新文丰出版公司据清光绪二十五年（1899）原刻本影印本。}
\BookNote{21.}{王钟翰，《清史补考·清代则例及其与政法关系之研究》，辽宁大学出版社，2004年，第73页。}
\BookNote{22.}{苏亦工，《明清律典与条例》，中国政法大学出版社，2000年，第74页。}
\end{BookNotes}

% 原书第 171 页（PDF 第 184 页）
\BookPage{171}{184}
\noindent 其衷也；四曰因援引尙涉狐疑，而申其议也；五曰酌量办理，而成式可循也；六曰系暂时变通，而非永着为例也；七曰虽经言官奏请，大吏条陈，因而酌改旧文，创立新例，尚未纂入条例也。」23这是就刑部通行章程而言，而各部门及地方勒定的章程，与现代的章程性质基本相同，在本部门及地方具有很强的司法效力，在一定时期和一定地域，发挥非常重要的作用。

成案，乃已成之案，是「国家之章程有定，而人事之变迁无穷，其有事出非常时，或为律所未及详者，则必议之群司，综诸宸断神明于律，而不谬于律，着为定例，以诏来兹」24。所以成案「俱系例无专条、援引比附加减定拟之案」25。在形式上，「成案是一种不成文的法律形式，是由各部或各省对某些典型案件判决的先例汇集而成的」。在本质上，成案，非正案，是根据律例而加减比附形成的案件判决。成案「至迟在明代已经出现了」26，清康熙年间就有「成案之多亦不啻数千」27。乾隆以后「随着司法权的下移，大量的成案被称引比附，并大有代例而行的趋向」28。虽然成案在此期间是否有代例而行的趋向还尚待例证，但是此段时间在司法审判大量征引比附确实是一种显著的现象2”。每一个案件都是一个综合复杂的情况，不可能通过简单的援引律例就可以分清责任，依法惩处。「今时律之外有例则已，备上下之比，而仍不能尽入于例，则又因案以生例，而其法详焉，故断狱尤视成案」。成案的成因是源于案情判决在律例之上无法寻找到合适的专条而加以征引，「其最善者莫如比照加减成案，事略而尽，文简而核，可以辅律例之未备」30。所以「成案是在律例没有规定之情况下适用的，目的是为了弥补成文法律的盲点」31。比如说在道光五年（1825），北京城内查获一起置买黄袍穿带的罪犯

\vfill
\begin{BookNotes}
\BookNote{23.}{王钟翰，《清史补考·清代则例及其与政法关系之研究》，辽宁大学出版社，2004年，第84页。}
\BookNote{24.}{（清）孙纶辑，《定例成案合镌·序》，康熙五十八年（1719）内府刊本。}
\BookNote{25.}{（清）祝庆祺，《刑案汇览凡例》，北京古籍出版社，2004年，第3页。}
\BookNote{26.}{苏亦工，《明清律典与条例》，中国政法大学出版社，2000年，第55页。}
\BookNote{27.}{（清）孙纶，《定例成案合镌》卷首〈序〉。}
\BookNote{28.}{（清)祝庆祺等编，《刑案汇览三编出版说明》，北京古籍出版社，2004年，第3页。}
\BookNote{29.}{参见《刑案汇览》、《刑部比照加减成案》、《刑部比照加减成案续编》、《新增刑案汇览》、《驳案汇编》、《驳案新编》等书。}
\BookNote{30.}{（清）许、熊莪，《刑部比照加减成案·熊莪叙》，清道光十四年（1834）刻本。}
\BookNote{31.}{何勤华，〈清代法律渊源考〉，载《中国社会科学》2001年第2期。}
\end{BookNotes}

% 原书第 172 页（PDF 第 185 页）
\BookPage{172}{185}
\noindent 张淳，这种行为则要依照谋大逆律凌迟处死，如果按照大逆罪就要籍没全家，但该人是「村野编氓」，所以「家属免其缘坐」32，显然是从情节上考虑。

此外中央部院衙门和地方官府颁布的禁约、告示、条约、告谕、檄文等可以统称为「禁约」。从管辖范围来看，其效力是不能够出本辖区，但一些具有普遍意义的禁约，往往会影响到其他地方。由于禁约是根据现实存在的问题，有针对性的制定出来的，既有约定俗成的惯例，也有申明律例的内容，有些经过整理汇编，可以成为省例，甚至经朝廷认可，成为固定章程，所以即便是部院衙门与地方府稍有过激的处置，在朝廷及上司的默许的情况下，也不属于违法。

由上可见，清代法规体系是以律、例、则例、事例、章程、成案等为主，以中央部院衙门和地方官府颁布的禁约为辅，构建了彼此相连的法规体系，因此在分析具体案件的时候，就要全面了解有关法规的逻辑结构，才能够深刻认识司法运行的动态和变化。

三

抄家案件除了反抗王朝统治的「谋反」、「叛逆」等重大事件所涉及的罪犯之外，其余案件大多都有政治方面的原因，既有复杂的人际关系，又有相互倾轧，还有君主的喜恶爱憎，同时也反映各级官僚贵族的一般财产和生活状况，还涉及到抄没财产的分配管理等问题，在具体处置上也存在比较大的差异。相同的罪行而不是相同的处置，从一些曾被抄家官员的处罚来看，明显区分为严惩、减等、从宽三种类型。

严惩是刑部等衙门会议拟罪以后，由皇帝直接核准，或皇帝在臣下拟罪以后加等的案件。这类案件不但查抄掉全部财产，而且还要籍没直系亲属。从批复的谕旨看，往往明确写着:「妻子俱充发某某地与披甲人为奴。」或「给驻防兵丁为奴」，甚至「给力能管束之回子为奴」、「充当折磨差使」，还有「妻子俱籍没入辛者库」。所谓「辛者库」，是满文“sinjeku”的音译，为内务府管辖的包衣组

\vfill
\begin{BookNotes}
\BookNote{32.}{（清）许楗，《刑部比照加减成案续编》卷5〈谋反大逆〉，清道光二十三年（1843）刊本。}
\end{BookNotes}

% 原书第 173 页（PDF 第 186 页）
\BookPage{173}{186}
\noindent 织，辛者库人为包衣管领下食口粮人，而被籍没的人口则为「缘罪入辛者库」，属于被监管的对象，原来是宗室贵胄或达官显宦及眷属，一经被拨入「辛者库」，便沦为比一般包衣门下更为低贱的身分。当然，被籍没入辛者库的人，都是旗人，汉人及其他民族则发配为奴或罚充苦差。

减等则是臣下拟罪较重，皇帝谕旨予以减轻一定的处罚。如臣下拟罪籍没家产人口，谕旨略给加恩，如「妻子免入辛者库」、「妻子免其为奴」等。这样即便是财产被抄没，眷属的身分不变，旗人可以继续领钱粮，非旗人可以参加科举考试，也算是「恩典」。

从宽是在臣下拟罪并要籍没家产人口时，皇帝从宽免于人口籍没为奴，并且给予一定的养瞻银两，当事人则从宽免死。如已故直隶总督杨景素被抄家时，乾隆帝在查抄奏折上亲笔批示：「着加恩将伊家产内酌量拨给三四万两俾资养瞻，其余分别估变解京」33。这是因为杨景素已死，所以格外开恩，当时查抄其财产有二十余万两，给还三四万两还可以让其子孙为富家翁。不过乾隆帝并没有满足，在他看来，「杨景素久历外任，贪黩营私，其原籍赀财，必不止仅有此数」34。直到三四年后，还继续清查其家族财产，并令其子承修河堤城垣工程。就是这样的「加恩」，皇帝也很少使用。

严惩、减等、从宽都取决于皇帝，这与皇帝的好恶有很大关系。以乾隆二十二年（1757）云贵总督恒文勒卖黄金勒索属员案来说，恒文与云南巡抚郭一裕都向皇帝进献金手炉，因向下属派购黄金的价钱不公道，以致全省沸腾。郭一裕惧怕事发牵连自己，便先发制人进行密奏，于是历时半年多的案件审理开始了。乾隆帝派钦差刘统勋前往查案，恒文供出是郭一裕出主意打造金手炉送给皇帝，并非仅本人压价购买黄金。问题是郭一裕不承认给皇帝送礼，而恒文却死称为了给皇帝送礼。在审理过程中，两人的任所及原籍家产都被查封，恒文财产估价37,466.877两，郭一裕财产估价40,809.65两。刘统勋依照律例拟恒文罪为绞监候，郭一裕杖一百、流三千里。然后由三法司核拟具奏，同意刘统勋的依律拟罪。接下来就要由皇帝做

\vfill
\begin{BookNotes}
\BookNote{33.}{《朱批奏折》内政类，〈职官〉，乾隆四十五年十月十三日，伊龄阿奏。}
\BookNote{34.}{《清高宗实录》卷1200，乾隆四十九年三月己亥条。}
\end{BookNotes}

% 原书第 174 页（PDF 第 187 页）
\BookPage{174}{187}
\noindent 出最后裁决。乾隆帝认为对恒文的处置太轻，「着派侍卫三泰、扎拉丰阿，驰驿前往，于解送所至之处，即将此旨宣谕，赐令自尽」。郭一裕则革职，「从宽发往军台効力，以为大吏鄙琐者戒」35。乾隆帝完全没有依照律例的规定，其中隐情，明眼人都能够看出，因为恒文「以进献为名」，将短价勒属的罪责推到皇帝身上；郭一裕不承认有「进献」，只是「惟以声色货利，殖产营运为事」，没有将责任推给皇帝，当然可以有「朕格外之恩」36。不死便有起复的机会，几年以后河南按察使空缺，乾隆帝「着郭一裕补授」37。

从本案审理过程，可以看到在错综复杂的案情，在各种人为因素面前，法律规定也变得模糊了。一般情况下，案件审理应该依照法律程式进行，但皇帝的态度决定案件审理的进程，皇帝喜怒及对全局的掌控又直接影响法律的适用。所谓雨露雷霆，恩威并济，一切都在皇帝掌控之中。如果将恒文案与李侍尧案比较，就可以发现，乾隆帝为了顾全自己的脸面，在惩处过程中根本不是什么折衷至当。「李侍尧身为大学士，历任总督，乃负恩婪索，盈千累万，甚至向屡员变卖珠子，赃私狼籍。如此不堪，实朕梦想所不到。不特朕用人颜面攸关，即各省督抚闻之，谅无不惭愧痛恨矣。」38其顾全的是情面，但李侍尧即使被判为斩刑，也没有讲出自己向乾隆帝进贡之事，而恒文却为了脱罪反复讲自己所为是因给皇帝进贡，两者孰生孰死，也就显而易见。史论云:「侍尧眷遇尤厚，屡坐赃败，屡屈法贷之。盖特怜其才，非以其工进献也。」39史称其才，而规避「进献」，实际上是清朝遗老遗少为「尊者讳」，如不进献，乾隆帝何能屡屈法贷之？李侍尧承认的赃私数量与抄没财产数量相比，抄没财产不及赃私的三分之一，其余三分之二都流向何方？不工进献才算怪呢？也就无怪乎乾隆帝惩贪鲜有成效了。

从涉及到政治经济等多种因素的重大抄家案件处置来看，按照律例规定，应该说政治因素高于经济因素，但在具体审判过程中却只谈经济因素，淡化政治因素，

\vfill
\begin{BookNotes}
\BookNote{35.}{《清高宗实录》卷546，乾隆二十二年九月辛丑条。}
\BookNote{36.}{中国第一历史档案馆编，《乾隆朝惩办贪污案件选编》，中华书局，1994年，第1册第59页。}
\BookNote{37.}{《清高宗实录》卷653，乾隆二十七年正月己未条。}
\BookNote{38.}{《清高宗实录》卷1103，乾隆四十五年三月丁酉条。}
\BookNote{39.}{《清史稿》卷323〈李侍尧传论〉，第10824页。}
\end{BookNotes}

% 原书第 175 页（PDF 第 188 页）
\BookPage{175}{188}
\noindent 多是探取「政治案件经济办」，既显示出法律为统治者工具性的特点，也表明政治斗争复杂性的特点。「不有废也，君何以兴？欲加之罪，其无辞乎？」（《左传·僖公十年》）君主权威低落，这是专制政体所不允许的，如果是权侔人主，必须将之除去，才能够确立君主的威严，也是帝王兴起之道，但如何处置，却是「君临之术」的重要内容。那么淡化政治犯罪而重惩经济犯罪，就成为君主及时了解和掌握政治全局，调整官制，予夺权势的重要手段之一了。比如说年羹尧有大逆、欺罔、贪渎等92款重罪，是震动全国的政治案件，雍正皇帝有许多批示，而于贪污而论，有贪黩之罪18款、侵蚀之罪15款，计有33款之多，是属于经济犯罪。政治罪是触犯皇权，完全可以依照君主的旨意将之定罪处死，但这样便会引到君臣关系上来，即便是君主绝对正确，在某些方面还会给人以皇帝不能够容人的感觉，因此明智的君主，从来不以政治罪来处置大臣。既然不以政治罪论处，其经济罪便会被提出来。年羹尧经济方面的犯罪就有33款之多，因此雍正帝仅以贪污来论罪。以经济罪论处，远比政治罪有说服力，因为其有贪污银两及物品是证据确凿，不容否认。这也是历代在处理权臣大案时，以政治罪定名少，以贪污罪定名多的原因之一。年羹尧及和坤案的处理则证实了上述论断。

君主专制中央集权制度实行的是奉上安下的体制，在奉上的环节常常会出现欺骗上、蒙蔽上、取悦上等变数，这一点在抄家案件中表现得尤为明显。从大量的抄家单来看，虽然是大如金银、土地、房产、商店、放债的借券，小到锅碗瓢盆、布衣碗筷之类，无不登录其内，有的抄家单厚达二三寸。但抄家单并不能反映出被抄者的真正财产，其中被抄者隐匿财产，查抄者欺隐财产，以及财产转移他人之手等情况也是屡见不鲜。例如，雍正查抄原直隶总督李维钧的抄家档案，最初抄家单仅有价值3,800两的财物，经过严查，则抄出现银34万余两、土地2,000多亩、房屋260余间，衣服用具等还没有计算价值。乾隆时查抄江苏巡抚庄有恭广东原籍，竟没有现银及房地产，要不是他病故而免于查抄，不可能没有现银及房地产的。据2000年3月《广州文物补查成果汇编》讲「庄有恭状元墓」的墓园广约200亩，因为有其父的坟墓，因此应该是庄有恭生前购置的，按抄家规定是可以计算入财产的。当然，这类督抚大员也有贪鄙清廉之分，财产也会有很大的差别，不过，被抄家的督抚们，很少能够称得上清廉。从督抚抄家单上看，他们大多有房屋200-


% 原书第 176 页（PDF 第 189 页）
\BookPage{176}{189}
\noindent 500间，土地480-3000亩，奴婢60-240名，其余金银首饰用器、绫罗绸缎衣服、字画古董等则难以估价。督抚都持有皇帝封赠自己及祖先诰命、御笔题字、御批文件、御赐图书用物等，这些在查抄时要「恭谨封固缴进」，不计算财产之内，但这些反映他们既往光荣和显赫身分的东西，对案件处理是否宽严有一定的作用。

这些高官一旦发现有人参劾，往往是转移财产，如曾任广东、广西巡抚的钱度被弹劾之后，派仆人往外运送财产，在贵州被查获值银五千以上的金玉诸器，在江西被查获银二万九千有奇，又从其原籍「书房地窖内起出银二万七千两，并寄顿金二千两，合计不下五六万两」40。在查获的过程中，也难保查获人不中饱私囊，或者以金换银。如原任浙江巡抚的王亶望，「诸州县馈赂率以千万计」，而被抄家时「得金银逾百万」41。所得与所存不符，后来发现浙江巡抚陈辉祖，在抄家时「有将金易银，抽换挪掩情事」42。由此可见，抄家单经过各级官吏之手，与实际抄没的数目应该存在很大差距。

抄没的财产与人口，按照律例规定，应该是入官和给主。官员因贪污受贿或亏空帑项，则所有的赃银赃物和追赔银两都要入官。抄没财产分现银与房地产什物两大类，房地产什物在拍卖以后，与现银一起计算数额。籍没人口分男女老幼，也有发配为奴、当差，以及先监禁后处置、老者赦免、妇女发卖等区别。因为涉及方面较多，当另行撰文予以剖析。

\vfill
\begin{BookNotes}
\BookNote{40.}{《清高宗实录》907，乾隆三十七年四月丁亥条。}
\BookNote{41.}{《清史稿》卷339〈王亶望传〉，第11074页。}
\BookNote{42.}{《清高宗实录》卷1164，乾隆四十七年九月甲辰条。}
\end{BookNotes}
